% 数理物理2 数理物理2-1 問題2（002-roots） 解いた日 2026-11-04
%% kotae
\toi{1}{◎}{}
\begin{align*}
-8+8\sqrt{3}\,\ii&=16\Bigl(-\frac{1}{2}+\frac{\sqrt{3}}{2}\ii\Bigr)\\
&=16\ee^{\ii\frac{2}{3}\pi}
\end{align*}
\toi{2}{○}{記述}
\begin{align*}
z^{4}&=16\ee^{\ii\frac{2}{3}\pi}\\
&=2^{4}\ee^{\ii\left(\frac{2\pi}{3}+\ayamari{\frac{k\pi}{2}}{2\pi k}\right)}\qquad k=0,1,2,3
\end{align*}
\naze{$z^{4}$ の偏角に足すのは $2\pi k$ で，$k\pi/2$ になるのは4で割った後の $z$ の偏角である。$z^{4}=2^{4}\ee^{\ii(2\pi/3+2\pi k)}$ から $z=2\ee^{\ii(\pi/6+k\pi/2)}$ と書く。$r$ は正の実数なので $r^{4}=16$ から $r=2$ だけをとること，$k=0,1,2,3$ で4個が出そろい $k$ を4変えると同じ解に戻ること，も書く。}
（図：原点を中心とする半径2の円と，円周上の4点を結んだ正方形）

半径2の円上に正方形
\[ z=2\ee^{-\ii\frac{5\pi}{6}},\ 2\ee^{-\ii\frac{\pi}{3}},\ 2\ee^{\ii\frac{\pi}{6}},\ 2\ee^{\ii\frac{2\pi}{3}} \]
\nuke{$a+b\ii$ の形}
\naze{問われている $a+b\ii$ の形にしていない。$-\sqrt{3}-\ii$，$1-\sqrt{3}\,\ii$，$\sqrt{3}+\ii$，$-1+\sqrt{3}\,\ii$ まで書く。}
\toi{3}{◎}{}
和 $0$　正方形の重心が原点だから。

積
\begin{align*}
&2^{4}\ee^{\ii\left(-\frac{5\pi}{6}-\frac{\pi}{3}+\frac{\pi}{6}+\frac{2\pi}{3}\right)}\\
&=16\ee^{-\ii\frac{\pi}{3}}=8-8\sqrt{3}\,\ii
\end{align*}
%% kizuki
\item $n$ 乗根は，右辺の偏角に $2\pi k$ を足してから $n$ で割る。$z^{n}$ の式に $2\pi k/n$ を書くと $z^{n}$ と $z$ が混ざる。$k=0,1,\ldots,n-1$ で $n$ 個が出そろい，$k$ を $n$ 変えると偏角が $2\pi$ 変わって同じ解に戻る。
\item 答えの形が指定されていたら（$a+b\ii$），その形まで直してから終える。最後に問題文の「〜の形で表せ」をもう一度読む。
